\documentclass{article}
\usepackage[utf8]{inputenc}
\usepackage[russian]{babel}
\usepackage{graphicx}
\usepackage{amsmath}
\usepackage{geometry}
\usepackage{booktabs}
\usepackage{pgfplots}
\usepackage{float}

\geometry{a4paper, margin=1in}

\title{Лабораторная работа №1:\\
Анализ чувствительности линейной регрессии}

\author{Тиганов В.И., ИСУ 467701, гр.J3112 \\ Вагин А.А., ИСУ 465339, гр.J3112 \\ Практический поток 1.3}
\date{12 марта 2025 г.}

\begin{document}

\maketitle

\section{Введение}
В данном отчете проводится анализ чувствительности модели линейной регрессии к ошибкам округления и исследуется сходимость вычисленных и исходных данных. Анализ основан на наборе данных \texttt{housing.csv}, с использованием различных конфигураций полиномиальных признаков и масштабирования. Цель работы — изучить влияние числа обусловленности матрицы признаков на точность модели и выявить оптимальные параметры для построения устойчивой регрессионной модели.

\section{Методология}
Набор данных разделен на обучающую и тестовую выборки с следующими параметрами:
\begin{itemize}
    \item Размер тестовой выборки: 0.23
    \item Параметр случайности: 45
\end{itemize}

К данным применяются полиномиальные признаки (степени 1, 2 и 3) и масштабирование (StandardScaler), после чего обучается модель линейной регрессии. Для оценки качества модели используются следующие метрики:
\begin{itemize}
    \item Средняя абсолютная ошибка (MAE) на обучающей и тестовой выборках
    \item Число обусловленности матрицы \( X_{train}^T \times X_{train} \) для оценки численной устойчивости
\end{itemize}

\section{Результаты}
Результаты экспериментов представлены в следующей таблице:

\begin{table}[H]
    \centering
    \begin{tabular}{cccccc}
    \toprule
    № эксперимента & Масштабирование & Степень полинома & MAE обучения & MAE теста & Число обусл. \\
    \midrule
    1 & Нет & 1 & 3.22 & 3.68 & \(7.28 \times 10^7\) \\
    2 & Нет & 2 & 1.80 & 2.56 & \(5.46 \times 10^{15}\) \\
    3 & Нет & 3 & 0.60 & 15.45 & \(8.48 \times 10^{30}\) \\
    4 & Да & 1 & 22.46 & 22.85 & \(9.08 \times 10^1\) \\
    5 & Да & 2 & 22.46 & 22.38 & \(2.11 \times 10^8\) \\
    6 & Да & 3 & 22.46 & 28.05 & \(2.60 \times 10^{18}\) \\
    \bottomrule
    \end{tabular}
    \caption{Результаты экспериментов с различными параметрами}
    \label{tab:results}
\end{table}

\section{Анализ}
\begin{figure}[H]
    \centering
    \includegraphics[width=0.9\textwidth]{graph.png}
    \caption{Зависимость числа обусловленности от MAE для обучающей (синий) и тестовой (оранжевый) выборок. Логарифмическая шкала по оси Y.}
    \label{fig:condition_vs_mae}
\end{figure}

Анализ результатов (Таблица \ref{tab:results} и Рисунок \ref{fig:condition_vs_mae}) позволяет сделать следующие выводы:

1. \textbf{Влияние масштабирования}:
\begin{itemize}
    \item Без масштабирования модель показывает лучшую точность на низких степенях полинома (MAE 3.22 против 22.46 при степени 1)
    \item Однако при увеличении степени полинома без масштабирования возникает проблема численной неустойчивости (число обусловленности достигает \(10^{30}\))
\end{itemize}

2. \textbf{Влияние степени полинома}:
\begin{itemize}
    \item Без масштабирования увеличение степени приводит к резкому росту числа обусловленности и переобучению (MAE обучения 0.60 vs MAE теста 15.45 при степени 3)
    \item С масштабированием ситуация более стабильна, но точность модели существенно ниже
\end{itemize}

3. \textbf{Оптимальная конфигурация}:
\begin{itemize}
    \item Наилучший баланс между точностью и устойчивостью достигается при степени полинома 2 без масштабирования (MAE теста 2.56)
    \item Однако при этом число обусловленности остается высоким (\(5.46 \times 10^{15}\)), что может вызывать проблемы при использовании разных реализаций алгоритмов
\end{itemize}

\section{Заключение}
Проведенный анализ позволяет сделать следующие выводы:

1. Масштабирование данных существенно снижает число обусловленности, но в данном случае приводит к значительному ухудшению точности модели. Это может быть связано с особенностями набора данных.

2. Увеличение степени полинома без масштабирования приводит к катастрофическому росту числа обусловленности и переобучению, несмотря на высокую точность на обучающей выборке.

3. Для данного набора данных оптимальным выбором является:
\begin{itemize}
    \item Использование полинома 2-й степени
    \item Отказ от масштабирования (в данном конкретном случае)
    \item Необходимость дополнительной регуляризации для улучшения устойчивости
\end{itemize}

4. Для промышленного применения рекомендуется:
\begin{itemize}
    \item Провести дополнительное исследование с другими методами масштабирования
    \item Рассмотреть возможность использования регуляризованных моделей (Ridge, Lasso)
    \item Провести анализ мультиколлинеарности признаков
\end{itemize}

Перспективы дальнейшего исследования включают анализ влияния различных методов регуляризации на устойчивость модели и точность прогнозов, а также сравнение с другими алгоритмами регрессии.

\end{document}